\origpage{14}{}
In the second case, we assume that $F$ separates $\Sigma$. Let $\Sigma_1, \Sigma_2 \subseteq \Sigma$ be surfaces with $\Sigma_1 \cup \Sigma_2 = \Sigma$ and $\Sigma_1 \cap \Sigma_2 = F$. Let $\Gamma_i = \Gamma \cap \Sigma_i$ $(i = 1, 2)$. Since $(\Gamma_1, \Gamma_2)$ is a separation of $\Gamma$ of order $< \theta' + 6 \le ord(\mathcal{T})$ (by (2)) we may assume from the symmetry that $(\Gamma_1, \Gamma_2) \in \mathcal{T}$. From (2), $\Sigma_1$ is not a disc, and so from (4.2)(iii), $\hat{\Sigma}_2$ is simpler that $\Sigma$. From (3.5) there is a tangle in $\Gamma_2$ of order $\ge ord(\mathcal{T}) - |F \cap V(\Gamma)| \ge f(f(\theta')) \ge f(\theta')$.

Now $\Gamma_2 \cap \Gamma'$ is obtained from $\Gamma_2$ by deleting some vertices and edges all incident with one region of $\Gamma_2$ in $\hat{\Sigma}_2$, and so from (1), $\Gamma_2 \cap \Gamma'$ has a tangle of order $\ge \theta'$. By (3.4), so does $\Gamma'$, as required. $\square$

By $\rho$ repeated applications of (4.5), we deduce

\textbf{4.6} \textit{For any surface $\Sigma$ with $bd(\Sigma) = \emptyset$, and all integers $\theta' \ge 1$ and $\rho \ge 0$, there exists $\theta \ge 1$ such that, if $\Gamma$ is a drawing in $\Sigma$ with a tangle of order $\ge \theta$, and $r_1, \dots, r_\rho$ are regions of $\Gamma$ in $\Sigma$, and $\Gamma'$ is obtained from $\Gamma$ by deleting some vertices and edges each incident with one of $r_1, \dots, r_\rho$, then $\Gamma'$ has a tangle of order $\ge \theta'$.}

\sectionhead{5 Linkages on surfaces}

In this section we prove (1.1) for graphs which can be drawn on a fixed surface. We need some further definitions.

If $C$ is a cuff of $\Sigma$ and $\Gamma$ is a drawing in $\Sigma$, there is a unique region of $\Gamma$ in $\hat{\Sigma}$ which includes $C \setminus V(\Gamma)$, and we call it the \textit{cuff region of $\Gamma$ in $\hat{\Sigma}$ corresponding to $C$}. An \textit{atom} of a drawing $\Gamma$ in a surface $\Sigma$ is either a region of $\Gamma$ in $\hat{\Sigma}$, or an edge of $\Gamma$, or a set $\{v\}$ where $v$ is a vertex. The set of atoms is denoted by $A(\Gamma)$, or $A_\Sigma(\Gamma)$ in cases of ambiguity. If $\Gamma$ is a drawing in $\Sigma$ and $\mathcal{T}$ is a respectful tangle in $\Gamma$, then, as discussed in [7], $\mathcal{T}$ defines a metric on $A(\Gamma)$, defined as follows. Let $K$ be a drawing in $\Gamma$ such that $V(\Gamma) \subseteq V(K)$, every region of $\Gamma$ includes a unique vertex of $K$, and $U(\Gamma) \cap U(K) = V(\Gamma) \cap V(K)$; and so that for every region $r$ of $\Gamma$, the vertex of $K$ it contains is adjacent in $K$ to every vertex $v$ of $\Gamma$ incident with $r$, by multiple edges if $r$ is incident with $v$ more than once, in the natural sense. For a closed walk $W$ of $K$ of length $< 2 ord(\mathcal{T})$, let $ins(W)$ be the union of the atoms of $K$ in $W$ together with all sets $ins(U(C))$ where $C$ is a circuit of $K$ whose edges all occur in $W$. For atoms $a, b$ of $\Gamma$, let $a', b'$ be the corresponding atoms of $K$ in the natural sense; we say $d(a, b) = 0$ if $a = b$, $d(a, b) = ord(\mathcal{T})$ if $a \ne b$ and there is no closed walk $W$ of length $< 2 ord(\mathcal{T})$ with $a', b' \in ins(W)$, and otherwise $d(a, b)$ is half the minimum length of such a walk. (See [7] for further discussion.) If $\{v\}$ is an atom of $\Gamma$ we often write $d(v, a)$ for $d(\{v\}, a)$. We call $d$ the \textit{metric of $\mathcal{T}$}.
